\documentclass[11pt,a4paper]{scrartcl}

\usepackage[T1]{fontenc}
\usepackage[utf8]{inputenc}
\usepackage[ngerman]{babel}
\usepackage{lmodern}
\usepackage{microtype}
\usepackage[a4paper,margin=2.5cm]{geometry}
% xcolor loaded below with table option
\usepackage{graphicx}
\usepackage{float}
\usepackage{setspace}
\usepackage{amsmath}
\usepackage{longtable}
\usepackage{multirow}
\usepackage{booktabs}
\usepackage{tabularx}
\usepackage{array}
\usepackage[table]{xcolor}
\usepackage[hidelinks]{hyperref}
\usepackage{xurl}

\definecolor{accent}{HTML}{0B4F6C}
\definecolor{colgreen}{HTML}{2E7D32}
\definecolor{colorange}{HTML}{E65100}
\definecolor{colred}{HTML}{C62828}
\definecolor{colblue}{HTML}{0B4F6C}
\definecolor{rowgray}{HTML}{F5F5F5}

\setstretch{1.15}
\setkomafont{disposition}{\color{accent}\bfseries}

\newcommand{\status}[1]{%
  \ifstrequal{#1}{Fertig}{\textcolor{colgreen}{\textbf{Fertig}}}{}%
  \ifstrequal{#1}{In Arbeit}{\textcolor{colorange}{\textbf{In Arbeit}}}{}%
  \ifstrequal{#1}{Geplant}{\textcolor{colblue}{\textbf{Geplant}}}{}%
  \ifstrequal{#1}{Weggelassen}{\textcolor{colred}{\textbf{Weggelassen}}}{}%
}

\begin{document}

% ───────────────────────────────────────────────────────
%  Titelseite
% ───────────────────────────────────────────────────────
\begin{titlepage}
  \centering
  \vspace*{3cm}

  {\color{accent}\rule{\linewidth}{2pt}}
  \vspace{0.8cm}

  {\LARGE\bfseries\color{accent} DAPI Zwischenstatusbericht\par}
  \vspace{0.4cm}
  {\large HydroNodeStation 01\par}
  {\normalsize Autarke Wetterstation inkl.\ KI-Analyse im HydroNode-Netzwerk\par}

  \vspace{0.8cm}
  {\color{accent}\rule{\linewidth}{1pt}}
  \vspace{2cm}

  \begin{tabular}{ll}
    \textbf{Autor:}       & Felix Knoll \\[4pt]
    \textbf{Kurs:}        & DAPI, Einzel-Projekt \\[4pt]
    \textbf{Datum:}       & 30.\ Mai 2026 \\
  \end{tabular}

  \vfill
  {\color{accent}\rule{\linewidth}{2pt}}
\end{titlepage}

% ───────────────────────────────────────────────────────
\section{Projektstatus}
% ───────────────────────────────────────────────────────

Das Projekt \textit{HydroNodeStation 01} befindet sich auf der Zielgeraden.
Der technische Kern ist abgeschlossen: Alle fünf Kernsensoren (Temperatur,
Luftfeuchtigkeit, Luftdruck, UV-Index, Feinstaub, CO${}_2$) liefern zuverlässige
Messwerte.
Die Station sendet Daten über LoRaWAN im Helium-Netzwerk; die Datenpipeline über
Amazon SNS zum HydroNode-Backend ist aktiv.
Website und iOS-App zeigen Echtzeit- und Verlaufsdaten an.

Der aktuelle Schwerpunkt liegt auf dem Hardwaretest im 3D-gedruckten
Stevenson-Screen-Gehäuse im Ausseneinsatz, die bisherigen Ergebnisse sehen gut aus.
Der TPS63900-Spannungsregler des PCBs hatte einen Fehler; als Workaround läuft
der Chip derzeit als externes Breakout-Board an den PCB-Headern, mit minimalem
Mehraufwand und gleicher Effizienz.
Das korrigierte PCB-Design ist bereits erarbeitet.
Offen sind noch die finale Solar-Verifikation, die KI-Evaluation und die
Abschlussdokumentation.


% ───────────────────────────────────────────────────────
\section{Soll-Zustand (Projektplan)}
% ───────────────────────────────────────────────────────

Im Projektplan wurden 20 Features mit insgesamt \textbf{160 Stunden} eingeplant
(Tabelle~\ref{tab:features}).

\begin{longtable}{>{\bfseries}p{0.5cm} p{13.5cm} r}
  \toprule
  \# & Feature / Aufgabe & Std. \\
  \midrule
  \endfirsthead
  \toprule
  \# & Feature / Aufgabe & Std. \\
  \midrule
  \endhead
  \bottomrule
  \endfoot
  \bottomrule
  \endlastfoot

  1  & Hardware-Auswahl \& Beschaffung (MCU, Sensoren, Solar, Akku, Gehäuse) & 8 \\
  \rowcolor{rowgray}
  2  & Schaltplan \& Aufbau der Wetterstation-Elektronik & 9 \\
  3  & Sensor: Temperatur, Luftfeuchtigkeit \& Luftdruck & 2 \\
  \rowcolor{rowgray}
  4  & Sensor: UV-Index \& Lichtstärke & 2 \\
  5  & Sensor: Partikelkonzentration (SPS30) & 4 \\
  \rowcolor{rowgray}
  6  & Sensor: CO${}_2$-Konzentration (SCD41) & 4 \\
  7  & Sensor: Umgebungslautstärke (MEMS-Mikrofon) & 3 \\
  \rowcolor{rowgray}
  8  & Standalone-Test: Sensoren \& LoRaWAN ohne Systemanbindung & 8 \\
  9  & PCB-Design, Bestellung \& Bestückung & 9 \\
  \rowcolor{rowgray}
  10 & Autarke Stromversorgung: Solar-Laderegler \& LiPo-Akku-Management & 13 \\
  11 & Deep-Sleep \& Power-Management Firmware & 7 \\
  \rowcolor{rowgray}
  12 & LoRaWAN-Firmware: Payload-Codec, Sendezyklus, Duty-Cycle-Management & 12 \\
  13 & Helium-Netzwerk: Registrierung, Console-Integration, Decoder & 5 \\
  \rowcolor{rowgray}
  14 & Backend: LoRaWAN-Datenpipeline zu HydroNode (Amazon SNS) & 8 \\
  15 & Website: Echtzeit-Dashboard für alle Stationssensoren & 14 \\
  \rowcolor{rowgray}
  16 & App: Historische Verlaufsdiagramme (alle Sensorkanäle) & 6 \\
  17 & App: Stations-Verwaltung (Hinzufügen, Benennen, Statusanzeige) & 4 \\
  \rowcolor{rowgray}
  18 & KI-Evaluation: Datensatz-Aufbereitung \& Modellauswahl & 8 \\
  19 & KI-Evaluation: Testing, Benchmarking \& Ergebnisdokumentation & 8 \\
  \rowcolor{rowgray}
  20 & Feldtest, Integrationstests \& Abschlussdokumentation & 16 \\
  \midrule
  \multicolumn{2}{r}{\textbf{Gesamt}} & \textbf{160} \\
\end{longtable}
\vspace{-1em}
\begin{center}
  \small\textit{Tabelle 1: Geplante Features mit Aufwandsschätzung aus dem Projektplan (Pitch, 26.03.2026)}
  \label{tab:features}
\end{center}


% ───────────────────────────────────────────────────────
\section{Ist-Zustand (aktueller Stand)}
% ───────────────────────────────────────────────────────

Bis zum 30.\ Mai 2026 wurden insgesamt \textbf{154 Stunden} in das Projekt investiert.
Der grösste Teil der geplanten Features ist abgeschlossen oder befindet sich in einem
fortgeschrittenen Stadium.

\subsection{Abgeschlossene Bereiche}

\begin{itemize}
  \item \textbf{Hardware \& Beschaffung:} Alle Komponenten wurden ausgewählt und bestellt.
        Als Hauptchip wird der STM32WLE5JC verwendet, der MCU und LoRa-Transceiver
        in einem einzigen Chip integriert und dadurch sehr energieeffizient ist.

  \item \textbf{Sensorik:} Folgende Sensoren sind in Betrieb und liefern zuverlässige
        Messwerte: SHT45 (Temperatur und Luftfeuchtigkeit), BMP390 (Luftdruck),
        LTR390 (UV-Index), SPS30 (Feinstaub) und SCD41 (CO${}_2$).
        Der im Plan vorgesehene Lautstärkesensor (MAX4466) wurde aus dem Scope
        genommen, da der Fokus auf den Kernsensoren liegt.

  \item \textbf{PCB-Design und Bestückung:} Ein eigenes PCB wurde in EasyEDA entworfen
        und über JLCPCB produziert. Das PCB ist bestückt und läuft stabil.
        Ein Fehler im TPS63900-Spannungsregler wurde identifiziert: Der Chip arbeitet
        nicht wie geplant. In Absprache mit einem Experten wurden Anpassungen erarbeitet;
        das korrigierte PCB-Design sollte für die bereits bestellten Boards passen.
        Als physischer Workaround wurde der fehlerhafte IC vom LiPo-Akku-Plus und dem
        3{,}3-V-Netz getrennt und der gleiche Chip als Breakout-Board an die PCB-Header
        angesteckt. Das System läuft damit identisch wie zuvor, mit minimalem
        Overhead und gleicher Effizienz.
        Das PCB-Design wird auf OSHWLab als Open-Source-Projekt veröffentlicht und
        nimmt am OSHWLab Open-Source-Contest teil.

  \item \textbf{Firmware:} Deep-Sleep und Power-Management sind implementiert und
        getestet. Der LoRaWAN-Stack mit Payload-Codec, Sendezyklus und
        Duty-Cycle-Management läuft stabil.

  \item \textbf{Helium-Netzwerk \& Backend:} Die Station sendet Daten erfolgreich
        über das Helium-Netzwerk. Die Datenpipeline über Amazon SNS zum HydroNode-
        Backend ist aktiv.

  \item \textbf{Website:} Ein Echtzeit-Dashboard für alle Stationssensoren wurde
        erstellt und ist öffentlich erreichbar. Die aktuellen Messwerte der Station
        können bereits live eingesehen werden unter:\\
        \url{https://hydronode.tech/stations/public/004dd7e7-5c27-4de0-bff5-24116f721658}

  \item \textbf{App-Erweiterungen:} Die iOS-App wurde um historische Verlaufsdiagramme
        für alle Sensorkanäle sowie um die Stationsverwaltung erweitert.

  \item \textbf{Gehäuse:} Das Gehäuse folgt dem Stevenson-Screen-Prinzip: übereinander
        angeordnete Lamellen halten Sonne und Regen ab, lassen aber Luft frei
        zirkulieren. Gedruckt in UV-beständigem ASA-Filament auf einem eigens
        angeschafften 3D-Drucker. Die Sensorik ist eingebaut; der Freilufttest läuft
        mit bisher guten Ergebnissen.
\end{itemize}

\subsection{Laufende Bereiche}

\begin{itemize}
  \item \textbf{Solarladung:} Das PCB läuft bereits mit Akku-Betrieb. Der TPS63900
        wird repariert, danach ist die vollständige Autarkie mit Solar gegeben.

  \item \textbf{KI-Evaluation:} Im HydroNode-Backend werden Sensordaten bereits
        automatisch auf Anomalien analysiert, auch mit den neuen Outdoor-Daten.
        Eine formale Evaluation mit Modellauswahl, Benchmarking und Dokumentation
        steht noch aus.

  \item \textbf{Feldtest \& Abschlussdokumentation:} Erste Langzeittests laufen
        bereits. Ein finaler Feldtest mit Aufstellung an der OTH Regensburg ist geplant.
\end{itemize}


% ───────────────────────────────────────────────────────
\section{Soll / Ist Vergleich}
% ───────────────────────────────────────────────────────

Tabelle~\ref{tab:vergleich} zeigt den Vergleich zwischen dem geplanten und dem
aktuellen Status für alle Features.

\begin{longtable}{>{\bfseries}p{0.5cm} p{6.5cm} r r l}
  \toprule
  \# & Feature & Geplant & Ist (est.) & Status \\
  \midrule
  \endfirsthead
  \toprule
  \# & Feature & Geplant & Ist (est.) & Status \\
  \midrule
  \endhead
  \bottomrule
  \endfoot
  \bottomrule
  \endlastfoot

  1  & Hardware-Auswahl \& Beschaffung & 8\,h & 7\,h & \status{Fertig} \\
  \rowcolor{rowgray}
  2  & Schaltplan \& Elektronik-Aufbau & 9\,h & 4\,h & \status{Fertig} \\
  3  & Sensor: Temperatur / Feuchte / Druck & 2\,h & 3\,h & \status{Fertig} \\
  \rowcolor{rowgray}
  4  & Sensor: UV-Index (LTR390 statt VEML6075) & 2\,h & 2\,h & \status{Fertig} \\
  5  & Sensor: Feinstaub SPS30 & 4\,h & 3\,h & \status{Fertig} \\
  \rowcolor{rowgray}
  6  & Sensor: CO${}_2$ SCD41 & 4\,h & 3\,h & \status{Fertig} \\
  7  & Sensor: Lautstärke (MAX4466) & 3\,h & 0\,h & \status{Weggelassen} \\
  \rowcolor{rowgray}
  8  & Standalone-Tests Sensoren \& LoRaWAN & 8\,h & 3\,h & \status{Fertig} \\
  9  & PCB-Design, Bestellung \& Bestückung & 9\,h & \textasciitilde56\,h & \status{In Arbeit} \\
  \rowcolor{rowgray}
  10 & Autarke Stromversorgung Solar \& Akku & 13\,h & \textasciitilde7\,h & \status{In Arbeit} \\
  11 & Deep-Sleep \& Power-Management & 7\,h & \textasciitilde10\,h & \status{Fertig} \\
  \rowcolor{rowgray}
  12 & LoRaWAN-Firmware & 12\,h & \textasciitilde11\,h & \status{Fertig} \\
  13 & Helium-Netzwerk Registrierung \& Decoder & 5\,h & \textasciitilde3\,h & \status{Fertig} \\
  \rowcolor{rowgray}
  14 & Backend: Pipeline zu HydroNode & 8\,h & \textasciitilde3\,h & \status{Fertig} \\
  15 & Website Echtzeit-Dashboard & 14\,h & 10\,h & \status{Fertig} \\
  \rowcolor{rowgray}
  16 & App: Verlaufsdiagramme & 6\,h & \textasciitilde6\,h & \status{Fertig} \\
  17 & App: Stationsverwaltung & 4\,h & \textasciitilde6\,h & \status{Fertig} \\
  \rowcolor{rowgray}
  18 & KI-Eval: Datensatz \& Modellauswahl & 8\,h & \textasciitilde2\,h & \status{In Arbeit} \\
  19 & KI-Eval: Testing \& Benchmarking & 8\,h & 0\,h & \status{Geplant} \\
  \rowcolor{rowgray}
  20 & Feldtest \& Abschlussdokumentation & 16\,h & \textasciitilde7\,h & \status{In Arbeit} \\
  \midrule
  \multicolumn{2}{l}{\textbf{Gesamt (inkl.\ 3D-Gehäuse, nicht geplant)}} & \textbf{160\,h} & \textbf{154\,h} & \\
\end{longtable}
\vspace{-1em}
\begin{center}
  \small\textit{Tabelle 2: Soll/Ist-Vergleich aller Features.}
  \label{tab:vergleich}
\end{center}

\noindent\textbf{Hinweis:} Die Ist-Stunden pro Feature sind Schätzungen. Im Arbeitstagebuch
wird nach Datum und Tätigkeit erfasst, nicht nach Feature-Nummer. Deshalb lassen sich
die Stunden nur näherungsweise zuordnen.


% ───────────────────────────────────────────────────────
\section{Stundenvergleich}
% ───────────────────────────────────────────────────────

Tabelle~\ref{tab:stunden} vergleicht die geplanten und tatsächlich aufgewendeten Stunden
gruppiert nach Bereich.

\begin{table}[H]
  \centering
  \begin{tabularx}{\textwidth}{X r r r}
    \toprule
    \textbf{Bereich} & \textbf{Geplant (h)} & \textbf{Ist (h, est.)} & \textbf{Differenz} \\
    \midrule
    Hardware, Beschaffung, Schaltplan & 17 & \textasciitilde9 & $-8$ \\
    Sensorik (5 von 6 Sensoren) & 12 & \textasciitilde11 & $-1$ \\
    Standalone-Tests & 8 & \textasciitilde3 & $-5$ \\
    PCB-Design, Bestellung, Inbetriebnahme & 9 & \textasciitilde50 & $+41$ \\
    Autarke Stromversorgung (Solar) & 13 & \textasciitilde7 & $-6$ \\
    Firmware (Deep-Sleep, LoRaWAN) & 19 & \textasciitilde21 & $+2$ \\
    Helium-Netzwerk \& Backend & 13 & \textasciitilde6 & $-7$ \\
    Website \& App & 24 & \textasciitilde23 & $-1$ \\
    KI-Evaluation & 16 & \textasciitilde2 & $-14$ \\
    Feldtest \& Dokumentation & 16 & \textasciitilde10 & $-6$ \\
    3D-Gehäuse (nicht geplant) & 0 & \textasciitilde12 & $+12$ \\
    \midrule
    \textbf{Gesamt (exakt laut Arbeitstagebuch)} & \textbf{160} & \textbf{154} & $-6$ \\
    \bottomrule
  \end{tabularx}
  \caption{Stundenvergleich Soll/Ist nach Bereichen (Ist-Werte geschätzt aus Arbeitstagebuch)}
  \label{tab:stunden}
\end{table}

\paragraph{Wichtigste Beobachtung:} Das PCB-Design hat mit ca.\ 50 Stunden weit mehr
Zeit gekostet als die geplanten 9 Stunden. Das entspricht einem Mehraufwand von rund
\textbf{+41 Stunden} in diesem Bereich. Dieser Mehraufwand wurde durch Einsparungen in
anderen Bereichen teilweise ausgeglichen: Helium und Backend liefen schneller als
erwartet, und die Standalone-Tests waren durch frühe Prototypentests bereits
abgedeckt. Der exakte Gesamtwert von 154 Stunden stammt direkt aus dem Arbeitstagebuch.


% ───────────────────────────────────────────────────────
\section{Abweichungen und Probleme}
% ───────────────────────────────────────────────────────

\subsection{PCB-Aufwand deutlich höher als geplant}

\textbf{Problem:} Im Projektplan waren 9 Stunden für das PCB eingeplant. Tatsächlich
wurden ca.\ 56 Stunden aufgewendet.

\textbf{Grund:} Das PCB wurde von Grund auf in EasyEDA neu entwickelt, da kein
geeignetes Open-Source-Design für den STM32WLE5JC mit allen benötigten Schnittstellen
existiert. EasyEDA ist ein browserbasiertes PCB-Design-Tool, das eng mit LCSC und
JLCPCB verknüpft ist und dadurch Bauteilauswahl und Fertigung vereinfacht. Hinzu kamen
mehrere Iterationen für RF-Design, Antennenmatching und Komponentenplatzierung. Nach
Erhalt der gefertigten Platine wurden ausserdem Bugs gefunden und behoben, unter
anderem ein fehlerhafter TPS63900-Spannungsregler. Das fertige Design wird auf
OSHWLab als Open-Source-Projekt veröffentlicht und nimmt am OSHWLab Open-Source-Contest
teil.

\textbf{Lösung:} Das Problem wurde in Absprache mit einem Experten analysiert; das
korrigierte PCB-Design soll für die bereits bestellten Boards funktionieren.
Als physischer Workaround wurde der fehlerhafte TPS63900 vom LiPo-Akku-Plus und dem
3{,}3-V-Netz getrennt und derselbe Chip als Breakout-Board an die PCB-Header angesteckt.
Das System läuft damit identisch wie zuvor, mit minimalem Overhead und gleicher Effizienz.
Die Solarladung wird nach finaler Verifikation freigeschaltet.

\subsection{Sensor-Wechsel}

\textbf{Situation:} Im Projektplan wurden bestimmte Sensoren als Kandidaten genannt.
Im Verlauf wurden leistungsfähigere Alternativen gewählt: BME280 wurde durch SHT45
und BMP390 ersetzt (genauer und energieeffizienter), VEML6075 durch LTR390
(breiterer UV-Messbereich, I2C statt analoger Ausgang).

\textbf{Bewertung:} Alle Sensoren laufen stabil. Die neuen Sensoren sind für den
Aussenbetrieb besser geeignet als die ursprünglich geplanten.

\subsection{Lautstärkesensor entfallen}

\textbf{Problem:} Feature 7 (Sensor Umgebungslautstärke, MAX4466, 3 Stunden geplant)
wurde aus dem Scope genommen.

\textbf{Grund:} Der Fokus lag auf den Kernfunktionen. Ein analoges MEMS-Mikrofon
erfordert zusätzliche Signalverarbeitung, die im verfügbaren Zeitrahmen nicht
sauber umgesetzt werden konnte.

\textbf{Lösung:} Das Feature wird als ``out of scope'' markiert. Die wichtigsten
Umweltdaten (Temperatur, Luftfeuchtigkeit, Luftdruck, UV, Feinstaub, CO${}_2$)
werden vollständig erfasst.

\subsection{KI-Evaluation noch nicht abgeschlossen}

\textbf{Problem:} Für die KI-Evaluation waren 16 Stunden geplant. Bisher wurden
erst ca.\ 2 Stunden investiert, da der PCB-Aufwand Priorität hatte.

\textbf{Lösung:} Die KI-Analyse ist bereits im Backend aktiv (Anomalieerkennung auf
Sensordaten). Eine formale Evaluation mit Modellauswahl und Dokumentation ist als
nächster Schritt geplant, sobald die Hardware vollständig stabil läuft.

\subsection{Nicht geplanter Mehraufwand: Gehäuse nach Stevenson-Screen-Prinzip}

\textbf{Situation:} Das Gehäuse-Design und der 3D-Druck wurde im Projektplan nicht
als eigenes Feature aufgeführt. Tatsächlich wurden ca.\ 11 Stunden für Design in
Fusion 360 und mehrere Druckdurchläufe aufgewendet.

\textbf{Warum dieser Aufwand nötig war:} Für eine korrekte Messung von
Aussentemperatur und Luftfeuchtigkeit darf sich das Gehäuse nicht durch Sonneneinstrahlung
aufheizen. Deshalb wurde das Design an einen Stevenson-Screen angelehnt: übereinander
liegende Lamellen blockieren direkte Sonneneinstrahlung und Regen, lassen aber
gleichzeitig Luft frei zirkulieren. Als Material wurde ASA gewählt, da ASA
UV-beständig, witterungsfest und für den dauerhaften Ausseneinsatz geeignet ist.
Der Druck erfolgte auf einem eigens für das Projekt angeschafften 3D-Drucker.

\subsection{Offener Punkt: UV-transparentes Glas für den UV-Sensor}

\textbf{Problem:} Der UV-Sensor LTR390 muss geschützt verbaut werden und trotzdem
UV-Licht empfangen können. Normales Glas oder transparentes PETG blockiert einen
grossen Teil der UV-Strahlung und verfälscht so die Messwerte.

\textbf{Stand:} Das passende UV-transparente Glas (Quarzglas oder Borosilikatglas)
würde mit Versand ca.\ 50 Euro kosten. Aktuell werden günstigere Alternativen
getestet, zum Beispiel dünne UV-transparente Kunststofffolien oder andere optische
Materialien, die UV-A und UV-B durchlassen. Dieses Detail ist für die Genauigkeit
der UV-Messung im Feldtest relevant.


% ───────────────────────────────────────────────────────
\section{Restarbeiten und Ausblick}
% ───────────────────────────────────────────────────────

Folgende Aufgaben stehen noch an:

\begin{enumerate}
  \item \textbf{Sensorplatine löten:} Eine eigene kleine Platine für den Stevenson-Screen
        wird noch gelötet. Sie wird über ein Datenkabel mit dem PCB verbunden und
        nimmt alle Sensoren passend auf.

  \item \textbf{3D-Druck Sensorhalterung:} Das abschliessende 3D-Modell für die
        Aufnahme der Sensorplatine im Stevenson-Screen muss noch gedruckt werden.

  \item \textbf{Solar-Autarkie:} Die vollständige Autarkie mit Solarversorgung muss
        noch validiert werden, insbesondere ob an bewölkten Tagen ausreichend Strom
        für den Dauerbetrieb erzeugt wird.

  \item \textbf{KI-Evaluation:} Datensatz aufbereiten, Modell auswählen
        und Ergebnisse dokumentieren.

  \item \textbf{Feldtest:} Station dauerhaft aufstellen (geplant an der OTH
        Regensburg), Langzeit-Integrationstests durchführen.

  \item \textbf{Abschlussdokumentation:} Alle Ergebnisse in der Studiendokumentation
        zusammenführen und finalisieren.
\end{enumerate}

\vspace{1em}
\noindent Der technische Kern des Projekts ist fertiggestellt. Die Station sendet
Sensordaten autark über LoRaWAN, die Daten sind in der App und auf der Website
einsehbar. Die verbleibenden Arbeiten betreffen hauptsächlich die Solar-Verifikation,
die KI-Evaluation und die Abschlussdokumentation.

\end{document}
